%% ma: L11-B17-0052
%% lop: 11
%% chuong: 5
%% bai: 17
%% dang: 11.17.4
%% loai: TN4
%% muc: TH
%% dap_an: "B"
%% nguon: "tai-lieu-hoc-tap-mon-toan-lop-11-theo-chuong-trinh-giao-duc-pho-thong-2018.pdf (Lớp Toán Thầy Duy, Phạm Lê Duy, THPT Chu Văn An), Chương 3 Bài 3, Phần 1, Câu 32 (tr.408)"
%% kien_thuc: "Lớp 11 Bài 17 (hàm số liên tục)"
%% trang_thai: cho-duyet
%% ly_do: "Dạng mới đề xuất (11.17.x), chờ giáo viên duyệt dạng; tài liệu gốc không có lời giải mẫu, đáp án do kho tự giải và kiểm bằng sympy."
\begin{ex}
Cho hàm số $y=\begin{cases}-x^2+x+3&\text{khi }x\ge2\\5x+2&\text{khi }x<2\end{cases}$. Chọn mệnh đề \textbf{sai} trong các mệnh đề sau.
\choice{Hàm số liên tục tại $x_0=1$.}{\True Hàm số liên tục trên $\mathbb R$.}{Hàm số liên tục trên các khoảng $(-\infty;2)$; $(2;+\infty)$.}{Hàm số gián đoạn tại $x_0=2$.}
\loigiai{\sloppy $f(2)=-4+2+3=1$, $\lim\limits_{x\to2^+}f(x)=1$, $\lim\limits_{x\to2^-}(5x+2)=12$. Hai giới hạn một bên khác nhau nên $f$ gián đoạn tại $x=2$, do đó mệnh đề "liên tục trên $\mathbb R$" sai. Các mệnh đề A, C (mỗi nhánh là đa thức), D đúng.}
\end{ex}
